\newpage
\section{PRESUPUESTO}

La viabilidad económica para el diseño, desarrollo y validación experimental del sistema comprende la adquisición de transductores, hardware embebido, placas de circuito impreso, consumibles de taller y servicios metrológicos. En la \tab{presupuesto_detallado} se presenta la estimación presupuestaria preliminar y tentativa desglosada por subsistemas técnicos en \textbf{Dólares Estadounidenses (\$USD)}, incorporando un margen de contingencia técnica y cambiaria del 15\%.

\begin{table}[H]
\centering
\caption{Presupuesto detallado para el desarrollo y validación del proyecto (en \$USD)}
\label{tab:presupuesto_detallado}
\vspace{1.0mm}
{\scriptsize
\setstretch{0.95}
\setlength{\tabcolsep}{2.8pt}
\renewcommand{\arraystretch}{0.90}
\begin{tabular}{|C{1.0cm}|L{7.4cm}|C{1.4cm}|R{1.8cm}|R{1.8cm}|}
\hline
\textbf{Ítem} & 
\textbf{Descripción del componente / servicio} & 
\textbf{Cant.} & 
\textbf{P. Unit. (\$)} & 
\textbf{Total (\$)} \\ \hline \hline

\multicolumn{5}{|l|}{\cellcolor{gray!12}\textbf{1. Subsistema de sensado e instrumentación analógica}} \\ \hline
1.1 & Sensor de corriente de núcleo partido no invasivo SCT013-000 (0--100~A) & 2 & 14.50 & 29.00 \\ \hline
1.2 & Sensor de corriente de efecto Hall integrado ACS712-20A (módulo calibrado) & 1 & 8.50 & 8.50 \\ \hline
1.3 & Acelerómetro digital triaxial MEMS ADXL345 (enlace SPI de bajo ruido) & 2 & 9.00 & 18.00 \\ \hline
1.4 & Componentes de acondicionamiento (filtros RC, zener 3.3V, resistencias 1\%) & 1 lote & 12.00 & 12.00 \\ \hline

\multicolumn{5}{|l|}{\cellcolor{gray!12}\textbf{2. Subsistema de procesamiento embebido, telemetría y potencia}} \\ \hline
2.1 & Placa de cómputo STM32F4 (ARM Cortex-M4 @ 84~MHz, FPU, CMSIS-DSP) & 1 & 24.00 & 24.00 \\ \hline
2.2 & Módulo microcontrolador Wi-Fi / Bluetooth ESP32 (nodo gateway MQTT) & 1 & 9.50 & 9.50 \\ \hline
2.3 & Fuente industrial conmutada 5V/2A con aislamiento galvánico y filtro EMI & 1 & 18.00 & 18.00 \\ \hline

\multicolumn{5}{|l|}{\cellcolor{gray!12}\textbf{3. Fabricación, encapsulado estanco y montaje magnético}} \\ \hline
3.1 & Fabricación de placa de circuito impreso PCB doble capa FR-4 industrial & 1 lote & 22.00 & 22.00 \\ \hline
3.2 & Gabinete IP65 y carcasa estanca IP67 con base magnética NdFeB N52 & 1 lote & 25.00 & 25.00 \\ \hline
3.3 & Cable balanceado apantallado y conectores industriales estancos IP67 (M8/M12) & 1 lote & 16.00 & 16.00 \\ \hline

\multicolumn{5}{|l|}{\cellcolor{gray!12}\textbf{4. Materiales fungibles de experimentación y mecanizado en MAINSA}} \\ \hline
4.1 & Barras cilíndricas de acero aleado AISI 1045 y AISI 4140 para ensayos de corte & 3 barras & 18.00 & 54.00 \\ \hline
4.2 & Caja de insertos intercambiables de carburo de tungsteno (ISO CNMG 120408) & 1 caja & 30.00 & 30.00 \\ \hline

\multicolumn{5}{|l|}{\cellcolor{gray!12}\textbf{5. Servicios metrológicos de calibración y laboratorio}} \\ \hline
5.1 & Rugosímetro Mitutoyo e inspección microscópica ISO 3685 & 1 serv. & 35.00 & 35.00 \\ \hline

\multicolumn{2}{|l|}{\textbf{SUBTOTAL DE RECURSOS DIRECTOS}} & \textbf{--} & \textbf{--} & \textbf{301.00} \\ \hline
\multicolumn{2}{|l|}{\textbf{CONTINGENCIA TÉCNICA Y CAMBIARIA (15\%)}} & \textbf{--} & \textbf{--} & \textbf{45.15} \\ \hline
\multicolumn{2}{|l|}{\textbf{COSTO TOTAL ESTIMADO DEL PROYECTO (\$USD)}} & \textbf{--} & \textbf{--} & \textbf{346.15} \\ \hline
\end{tabular}
}
\vspace{1.5mm}
\small \textbf{Fuente:} Elaboración propia, 2026.
\end{table}

\vspace{0.15cm}
El costo total estimado para la ejecución experimental del proyecto asciende a \textbf{\$346.15 USD}, contemplando una reserva del 15\% (\$45.15 USD) frente a contingencias y variaciones de cambio. La totalidad de los recursos será asumida por el postulante, mientras que MAINSA facilitará horas de torno CNC, herramientas de sujeción, fluidos refrigerantes y probetas de trabajo. Es indispensable precisar que los costos de la \tab{presupuesto_detallado} poseen un \textbf{carácter preliminar y tentativo}, formulado estrictamente para el prototipo de laboratorio y banco de pruebas; durante la fase de ejecución en Taller de Grado II, estos rubros se contemplarán y reajustarán con mayor cautela analítica mediante cotizaciones comerciales formales.

En términos de viabilidad económica, los sistemas comerciales propietarios de monitoreo de herramientas (TCM) en el mercado internacional, basados en dinamómetros de cuarzo o consolas cerradas, demandan entre \$10,000 y \$30,000 USD, resultando prohibitivos para maestranzas bolivianas. El sistema propuesto representa menos del 3.5\% de dicha inversión, demostrando alta viabilidad, rápido retorno y un relevante valor de transferencia tecnológica.

Asimismo, se establece una clara distinción metodológica entre la fase de prototipado de laboratorio y la especificación de grado industrial. Para la validación experimental se seleccionaron componentes de alta disponibilidad local y bajo costo (transductores SCT-013, microcontrolador STM32F4 y acelerómetros MEMS), complementados con encapsulados estancos IP67, acoplamiento magnético NdFeB N52, aislamiento galvánico de señales analógicas y filtros EMC. Esta estrategia permite validar la correlación física y el desempeño algorítmico a coste mínimo, sentando las bases para su posterior escalamiento a circuitos SMD industriales.

Atendiendo a las recomendaciones del comité evaluador sobre robustez en planta, en el \textbf{Anexo 2} se presenta de forma complementaria la estimación presupuestaria proyectada con componentes de grado industrial robusto (transductores certificados 4--20~mA, acelerómetro piezoeléctrico IEPE/IP68 en acero inoxidable, acondicionamiento de riel DIN y envolventes IP66). Dicho análisis ratifica que, aun con componentes de máxima exigencia industrial, la solución proyectada conserva un ahorro superior al 88\% frente a los sistemas comerciales de importación.


